%!TEX root = ../hutbthesis_main.tex

% 设置中文摘要
\keywordscn{AirSim，无人机集群，仿真平台，性能优化，线程调度}
\begin{abstractzh}

针对微软AirSim平台在大规模无人机集群仿真中帧率骤降、时延激增的瓶颈，本文围绕无人机集群仿真、性能优化及线程调度展开软件层优化研究。低效线程调度、内核态与用户态频繁切换、不合理自旋等待是导致性能退化的主要原因。单线程分发器、标准阻塞睡眠、协作式让出等五种线程等待模型被设计并集成于AirSim框架。微观基准测试与集成仿真实验从定时精度、帧率、CPU占用等维度验证了优化效果。经微观基准测试与集成仿真实验，从定时精度、帧率、CPU占用等维度完成量化评估。结果表明，优化方案显著提升AirSim支持20架以上无人机仿真的效率与稳定性，为高并发实时仿真提供可复用调优范式，降低集群算法真机测试成本，推动无人机集群在侦察、配送、测绘等领域落地应用。

\end{abstractzh}
